\chapter{What Happens in Solutions}

\begin{kaichang}
Try a small kitchen action: fill a glass with clean water, add a spoonful of salt, and stir. Ten seconds later, the salt is gone. The glass still holds transparent liquid without visible particles---but taste it, and it is salty.

Make guesses before reading the answers to three questions:
\begin{itemize}[itemsep=2pt]
\item Did the glass become heavier when the salt dissolved, or did salt ``lose'' its weight?
\item Are the top and bottom equally salty? What about after a month?
\item If you leave the water to evaporate, will the salt return?
\end{itemize}

Write down your guesses. By the end, you will see that ``the salt disappeared'' may be the kitchen's most misleading statement. At particle level, nothing vanished; an astonishingly large rearrangement occurred.
\end{kaichang}

\section{Weigh First: Nothing Disappears}

Before explanations, do what any scientist would do first: weigh.

Place \ud{250}{g} of water on a digital scale and record the reading. Add \ud{5}{g} of salt and stir. The result? \ud{255}{g}, with nothing missing. Seal the bottle and weigh it a month later: still \ud{255}{g}. The salt lost no weight; the scale simply cannot ``see'' where it is now.

Observe a few more facts:

\begin{itemize}[itemsep=2pt]
\item Stirred saltwater tastes equally salty at the top and bottom, however long it stands. Salt does not settle like sand. Stir sand into water and layers form in minutes; saltwater remains uniform for a year.
\item Pour saltwater into a shallow dish and let it evaporate slowly. White crystals return. Taste them: salt. Weigh them: close to the original \ud{5}{g}.
\item Salt is not alone: dissolved sugar tastes sweet, a cold-remedy granule drink flavors the whole cup, and soda contains invisible yet perceptible carbon dioxide. Iodine tincture dissolves iodine in alcohol. \textbf{The solvent need not be water}, though water is most common.
\end{itemize}

Uniform solute distribution and no separation or settling however long it stands: these define a \textbf{solution}. The dissolved material is the \textbf{solute} (salt, sugar, carbon dioxide); the dissolving medium is the \textbf{solvent} (water, alcohol).

Thus dissolving really means \textbf{separating and dispersing solute particles uniformly among solvent particles}. They remain, just too small to see. How small? Time to zoom in.

\section{Zooming In: Water Dismantles the Walls}

Recall two earlier accounts. Chapter 18 showed that salt contains no independent ``\chem{NaCl} molecules,'' but a huge checkerboard crystal of \chem{Na^+} and \chem{Cl^-} held by opposite-charge attraction. Chapter 23 showed bent water with a negative oxygen end and positive hydrogen ends, a tiny ``electromagnet.''

Put salt in water and an invisible siege begins. Exposed \chem{Na^+} is immediately surrounded by water's oxygen ends; surface \chem{Cl^-} by hydrogen ends. Dozens of molecules tug an ion loose, then keep it wrapped in successive layers like a ``water shell.'' This surrounding process is \textbf{solvation}, specifically \textbf{hydration} when water is the solvent. Constantly jostled hydrated ions wander randomly and diffuse from bottom to surface, explaining why every sip tastes the same.

Sugar takes another route. Electrically neutral sucrose cannot split into ions, so water surrounds and carries away whole molecules. They float intact, keeping sugar water sweet: sweetness belongs to sucrose, and an undamaged molecule keeps its taste.

One phrase matters: \textbf{always moving}. Dissolution does not end with putting salt into water; ions and molecules join the crowd and collide ceaselessly day and night. A quiet-looking glass contains an endless crowded dance.

How large a dance? Salt's molar mass is about \ud{58.5}{g/mol}, so \ud{1}{g} is $1/58.5\approx 0.017$ mol. A mole contains $\ud{6.02\times 10^{23}}{}$ ``\chem{NaCl} units,'' each splitting into two ions. Thus \ud{1}{g} produces about $0.017\times 6.02\times 10^{23}\times 2\approx 2\times 10^{22}$ ions. A small spoonful in one glass averages around $10^{17}$ ions per cubic millimeter, far exceeding all the sand grains on Earth's beaches. Your eyes cannot cope with particles so tiny; no wonder they seem gone.

\section{Limits to Dissolving: Saturation, Supersaturation, and Bubbles}

If water molecules are movers, do they have a workload limit? Yes, as you know from repeatedly adding and stirring salt. Eventually new salt sits on the bottom and will not diminish. The solution is \textbf{saturated}.

But saturation does not mean stopping. Water still removes surface ions, while wandering hydrated ions collide with the crystal and rejoin its lattice. \textbf{Dissolving and crystallizing happen simultaneously}. Initially departures exceed returns, so salt decreases. With enough dissolved ions, returns catch up, exactly balancing departures. The calm surface hides opposing particle floods. Such dynamic equilibrium is everywhere in chemistry and receives dedicated attention in Chapter 28.

A classic piece of evidence: suspend an irregular salt grain by thread in saturated saltwater for several days. Its weight barely changes because flows balance, but its edges smooth as surface particles are exchanged. If saturation meant nothing happened, it would remain untouched.

The saturation threshold follows two controls:

\begin{itemize}[itemsep=2pt]
\item \textbf{Temperature}. Most solids dissolve more in hotter water. Potassium nitrate, \chem{KNO_3}, dissolves about \ud{32}{g} per \ud{100}{g} water at $20\,^{\circ}\mathrm{C}$, but \ud{169}{g} at $80\,^{\circ}\mathrm{C}$, over fivefold more. Cool a hot, fully loaded solution and excess solute must surrender its place as crystals. Sugar grains in stored honey and rock sugar growing in syrup use this principle.
\item \textbf{Pressure}. Solids and liquids are barely affected, but for \textbf{gases} pressure matters greatly. Soda packs \chem{CO_2} into water under high pressure. Open the cap, pressure plunges, and excess gas rushes out as bubbles.
\end{itemize}

Temperature works oppositely for gases: hotter water dissolves \textbf{less}. Fish gasp at pond surfaces on sultry summer afternoons because warm water lacks oxygen. Bubbles on a cold-water glass left at room temperature are dissolved air ``checking out.''

There is also a tricky state. Carefully cool a hot concentrated solution without vibration, and solute may fail to crystallize promptly, leaving it ``overstaffed but temporarily uninspected'': \textbf{supersaturation}. This is highly unstable. Add a seed crystal, or sometimes merely shake it, and excess solute crystallizes in seconds into a forest. Reusable winter hand warmers triggered by flexing a metal disc contain supersaturated sodium acetate solution. Chapter 27's energy ledger explains the heat released during crystallization.

\section{Sugar Water Cannot Light a Bulb; Saltwater Can}

One comparison holds the key to solution conductivity.

Imagine a battery-and-bulb circuit with exposed metal electrodes at two wire ends. Insert them into purified water: almost no light. Concentrated sugar water: still none. Saltwater lights the bulb; white vinegar gives a faint glow.

Current is directed movement of charged particles. In solid metals, electrons carry it; liquids lack free electrons and need \textbf{freely moving charged particles}. Whole water and sugar molecules are neutral, so sugar water has nobody to carry charge. Saltwater contains abundant wandering \chem{Na^+} and \chem{Cl^-}. Connect power and positives head for the negative electrode, negatives for the positive, relaying current through two ion streams.

Substances conducting when dissolved in water or molten are \textbf{electrolytes}; those that do not are \textbf{nonelectrolytes}. Most salts, acids, and bases are electrolytes; sugar and alcohol are not. Recall that solid salt does not conduct because its ions are locked in place (Chapter 18). Same salt, different conductivity, solely because ion mobility changes.

Now symbols can enter. We use (s) for solid and (aq) for the hydrated state, from Latin aqua, water:

\[\chem{NaCl(s)} \rightarrow \chem{Na^+(aq)} + \chem{Cl^-(aq)}\]

One \chem{NaCl} unit becomes one hydrated sodium and one hydrated chloride ion. Sucrose's equation is simpler:

\[\chem{C_{12}H_{22}O_{11}(s)} \rightarrow \chem{C_{12}H_{22}O_{11}(aq)}\]

Its molecule stays intact, moving from solid into water. Side by side, these equations explain conductivity: both start as molecules or lattices, but only one yields charged particles.

This is more than a classroom trick. Sweat is a genuine electrolyte solution, making wet-hand contact with switches much more dangerous; ions in the water film open a path for current. Hospital saline, aquarium water changes, and boiler-scale treatment all revolve around which ions water contains and how many.

\section{How Scientists Know Salt Really Splits in Water}

\begin{zhengju}
Today, dissolved salt becoming ions seems obvious, but in the 1880s it seemed absurd. The prevailing view held that current split compounds into ions: electricity was cause, ions effect.

Trouble arose with van 't Hoff. Around 1885, he systematically measured solutions' \textbf{osmotic pressure}, discussed shortly; for now, imagine their ``thirst for water.'' Sugar solutions perfectly matched a formula giving higher pressure for more solute particles, as though they behaved independently like gas molecules. Salt at the same concentration gave almost \textbf{twice} the predicted pressure, while \chem{CaCl_2} approached \textbf{three times}. Why different rules for salt and sugar?

In his 1884 doctoral thesis, the young Swedish student Arrhenius proposed that electrolytes \textbf{automatically} dissociated into charged ions upon entering water, without current. One \chem{NaCl} yielding two particles naturally doubles osmotic pressure; \chem{CaCl_2} yielding three triples it. Van 't Hoff's odd data immediately made sense. The ratio became the van 't Hoff factor.

Professors rejected the idea, awarding only a barely passing thesis grade. Their objection sounded strong: sodium explodes with water and chlorine is highly poisonous, yet sodium and chloride ions supposedly float peacefully in it? Arrhenius's simple, profound reply was that \textbf{ions and atoms are different things}. Sodium atoms are reactive because they have an electron to lose; ions have already lost it and are mild-mannered, as Chapter 10's shell logic explains. Dissolved \chem{Na^+} versus sodium metal resembles a retiree versus a boxer: the same person in utterly different states.

Debate continued until independent evidence converged: conductivity, freezing-point changes, and X-ray images of ionic crystals all agreed. Arrhenius received the 1903 Chemistry Nobel. Then came another chapter: concentrated solutions had smaller factors, with \chem{NaCl}'s doubling falling to 1.9 or 1.8. Dissociation theory was not wrong; crowded ions interfered and acted together. Debye and H\"uckel calculated this correction in 1923. Scientific theories rarely settle everything at one stroke; they gain reinforcement layer by layer.
\end{zhengju}

\section{Between Dissolved and Undissolved: Colloids and Milk}

Now challenge our categories. Milk looks uniform and does not settle for days. Is it a solution?

Interrogate it with light. In darkness, direct a laser pointer or phone torch through saltwater and diluted milk. The light path is almost invisible in saltwater but bright in milk. This phenomenon, discovered by the British physicist Tyndall, reveals \textbf{particles large enough to scatter light}. Saltwater ions measure below a nanometer and barely affect the passing waves; milk's fat globules and protein aggregates are hundreds of nanometers across and scatter them.

Particle size arranges dispersions along a continuum:

\begin{itemize}[itemsep=2pt]
\item \textbf{True solutions} (below about 1 nanometer): ions and small molecules, transparent and never settling, such as saltwater and sugar water.
\item \textbf{Colloids} (about 1 nanometer to 1 micrometer): small enough not to settle but large enough to scatter light, including milk, soy milk, blood, jelly, fog, and clouds.
\item \textbf{Suspensions and emulsions} (above about 1 micrometer): eventually settle or separate, such as muddy water and poorly mixed oil and water.
\end{itemize}

Milk is also a stabilized \textbf{emulsion}. Oil and water normally separate quickly because polar and nonpolar molecules avoid one another (Chapter 23). Milk proteins have water-loving and oil-loving ends, wrapping droplets with hydrophilic ends outward, disguising oil as something water welcomes and stabilizing its suspension. Such substances are \textbf{emulsifiers}. Egg-yolk lecithin emulsifies mayonnaise; detergent similarly pries grease from dishes, wraps it in water, and washes it away.

Colloids are not obscure textbook terms: blood plasma, tofu setting, blue skies, and water-treatment flocculation all belong to colloid chemistry.

\section{A Membrane That Admits Some Particles: Osmosis}

Our last stop is the chapter's most important concept, explaining a pickled cucumber today and every cell tomorrow.

Some membranes admit small water molecules but exclude larger solute particles. These are \textbf{semipermeable membranes}: the white membrane inside eggshells, sausage casings, and laboratory cellophane are examples.

Imagine dividing a tank with such a membrane, pure water on one side and concentrated sugar solution on the other. Hours later, sugar solution rises and pure water falls. \textbf{Water crosses the membrane with a net flow toward sugar solution}. This is \textbf{osmosis}.

Why the concentrated side? Reject a tempting explanation: concentrated solution does not extend a hand and suck water across. The truth is statistical. Water on both sides moves randomly and may cross pores, but solute occupies some membrane-contact positions on the concentrated side, reducing water crossings from that side. Both directions continue; more comes from dilute than concentrated solution, producing net flow toward concentration. Eventually the raised liquid column supplies enough pressure to push the extra water back and rebalance flows. This equilibrium pressure difference is \textbf{osmotic pressure}.

\begin{qiaoliang}
How large is it? Van 't Hoff found an almost gas-like formula, $\pi = iMRT$, where $M$ is molar concentration, $i$ the number of particles produced per solute unit, $R$ a constant, and $T$ absolute temperature. Consider hospital saline, $0.9\%$ salt solution: \ud{9}{g} per liter gives $9/58.5\approx 0.154$ mol. Two ions per unit give $i\approx 2$. With $T\approx 310$ K (body temperature) and $R\approx \ud{0.082}{L\cdot atm/(mol\cdot K)}$:

\[\pi \approx 2\times 0.154\times 0.082\times 310 \approx 7.8\ \mathrm{atm}\]

About 7.8 atmospheres, equivalent to an 80-meter water column! Half a small spoonful in a bottle mobilizes enormous force, so osmotic pressure deserves respect. Those 7.8 atmospheres exactly match body fluids, explaining $0.9\%$ infusion saline: \textbf{isotonic}, neither too much nor too little. Pure water directly in a vein makes red cells absorb water and burst; concentrated saltwater makes them shrink. Small dose differences put cells' lives on the line.
\end{qiaoliang}

Osmosis pervades life. Salt draws cucumber-cell water into pickling juice; sugared tomatoes release the same sweet puddle. Marine fish drink seawater and excrete salt through gills, while freshwater fish scarcely drink and urinate profusely. Each evolved its own osmotic strategy, and exchanging environments is fatal.

Humans reverse the process: pressure \textbf{exceeding} seawater's osmotic pressure forces water backward through a membrane, leaving salt behind. This is \textbf{reverse-osmosis} desalination. Much drinking water in Israel, Singapore, and Persian Gulf cities is squeezed from the sea this way. Your purifier probably has a reverse-osmosis membrane cartridge too.

\begin{shiyan}
\textbf{``Naked Eggs'' That Gain and Lose Water}

Materials: 2 raw eggs, a bottle of white vinegar, concentrated sugar solution or honey, clean water, 3 transparent glasses, cling film, and optionally a kitchen scale.

Procedure:
\begin{enumerate}[itemsep=2pt]
\item Gently place eggs in a glass, cover with vinegar, and cover the glass with film. Tiny bubbles immediately appear as acetic acid dissolves shell calcium carbonate and releases \chem{CO_2}.
\item Leave 24--48 hours, replacing vinegar once. When shells disappear, translucent elastic naked eggs remain inside their natural semipermeable membrane, admitting water but not sugar or proteins.
\item Put egg A in concentrated sugar solution and egg B in clean water for 12 hours, turning them if desired.
\end{enumerate}

What to observe: A wrinkles, shrinks, and loses mass (weigh before and after if possible); B swells, gains mass, and feels taut. Egg fluid is more concentrated than water but less than sugar solution, so net water flow always favors the more concentrated side.

Safety note: vinegar's smell is irritating; ventilate and rinse splashed eyes immediately with plenty of clean water. Eggs are contaminated by vinegar and outside microorganisms and \textbf{must not be eaten}. Handle glasses gently.
\end{shiyan}

\section{Where This Picture Fails}

Our model is neat and useful, but know its boundaries before it misleads you:

\begin{itemize}[itemsep=2pt]
\item Higher concentration giving higher osmotic pressure, and van 't Hoff factors equaling ion counts, hold precisely only for \textbf{dilute solutions}. Concentrated ions interfere and water becomes scarce, departing from simple formulas. Such corrections helped establish Arrhenius's theory.
\item Increasing solubility with temperature fits most solids but not calcium hydroxide, slaked lime, which becomes less soluble, while gases generally do the opposite. Empirical patterns are not laws; check data first.
\item Dissolving and reacting have no sharp wall between them. Salt's ions remain unchanged, so we call dissolution physical; some dissolved \chem{CO_2} becomes carbonic acid, while sodium metal reacts explosively. Pure mixing and complete transformation lie along a continuum. The next part, beginning in Chapter 26, handles transformation formally.
\item The boundary at 1 nanometer between colloids and true solutions is conventional. Nanoparticles, giant proteins, and viruses straddle it. Categories are tools, not cracks in nature.
\end{itemize}

\begin{wuqu}
``Osmosis is water attracted to high concentration.'' Personification causes this error. Molecules cannot see concentration across the membrane, and no suction spans it. Each merely collides randomly, wherever it happens to go. Net direction comes from statistics: solute replaces some water molecules approaching the membrane on the concentrated side. A simple test turns the U-tube sideways or upside down: osmosis continues independently of gravity. Attraction language cannot explain this; statistics can. Remember our refrain: ask not where particles want to go, but which route more of them take.
\end{wuqu}

\section{Where Did That Spoonful Go?}

Check our three opening guesses.

First, saltwater becomes heavier by the full combined weights. Mass conservation rests on atom conservation: dissolution neither creates nor destroys atoms, only reorganizes an ordered lattice into wandering hydrated-ion armies. Chapter 26 establishes the universal rule behind it.

Second, every sip is equally salty and always remains so. About $2\times 10^{22}$ ions wander uniformly through thermal motion, neither settling nor layering nor returning to the bottom.

Third, evaporate the water and ions lose their watery surroundings, reassembling into the alternating checkerboard. The same salt returns.

Disappearance is not a particle's fate, only your eyesight's limit.

\section{From Saltwater to a Membrane-Bound City}

This chapter has quietly prepared a bigger story.

Blood, sweat, tears, and cytoplasm are solutions. An adult body is a precisely regulated aqueous system of about 40 liters. Dissolved glucose travels through blood, sodium and potassium generate nerve signals, and kidneys continually regulate your ``glass of solution.''

All this requires every cell's surrounding membrane, first of all semipermeable: water passes, many solutes cannot pass at will. Our questions become matters of cellular life and death. Uncontrolled osmotic pressure bursts or shrivels cells; ionic imbalance stops nerves and hearts. But cell membranes exceed eggshell membranes in sophistication, adding one-way ``water pumps'' and ion-specific gates to turn passive osmosis into active management. How does a nonliving membrane make seemingly thoughtful choices? Volume Two's Chapter 51, ``Cell Membranes,'' answers. Hydrated ions, concentration, semipermeability, and osmotic pressure all return unchanged, on a stage that shifts from beaker to you.

\begin{tiaozhan}
Solute particles also change solvent properties by count alone, irrespective of identity: colligative properties. Besides osmosis, these include \textbf{freezing-point depression} and \textbf{boiling-point elevation}. In dilute solutions, $\Delta T = i\cdot K_f\cdot m$, where $K_f$ is the solvent constant (water: $1.86\,^{\circ}\mathrm{C\cdot kg/mol}$) and $m$ molality. For winter, add 1 mol salt (\ud{58.5}{g}) to \ud{1}{kg} of melted snow water. With $i\approx 2$, freezing drops about $2\times 1.86\approx 3.7\,^{\circ}\mathrm{C}$, explaining road salting. But below roughly $-10\,^{\circ}\mathrm{C}$, no amount helps, requiring calcium chloride ($i\approx 3$, with heat release) or snowplows. Reverse the application: sweet ice-lolly mixtures freeze less readily than water. Test whether your freezer can freeze saturated saltwater solid.
\end{tiaozhan}

\section*{Try It}

\begin{sikaoti}
\item Draw a saltwater drop magnified a hundred million times, showing \chem{Na^+}, \chem{Cl^-}, and surrounding water with negative oxygen and positive hydrogen ends correctly oriented. Note that this is a representation; hydration shells continually disassemble and reform.
\item Suspend rough rock sugar in its own saturated solution, seal, and leave a week. Predict mass and shape changes through opposing dissolution and crystallization. Why would this demonstrate that saturation is not stillness?
\item Draw flows across a membrane separating pure water on the left and concentrated saltwater on the right. Use arrow thickness for flow amounts, mark net direction and final equilibrium, then replace pure water with dilute saltwater. Does direction change?
\item Why do doctors give dehydrated patients $0.9\%$ saline rather than pure water or $9\%$ saline? Treat red cells as water-filled balls in semipermeable membranes and draw consequences of both wrong choices.
\item Given clean water, diluted milk, saltwater, and a laser pointer, design an experiment distinguishing all three and predict observations.
\item Estimate moles of ions and ion count in \ud{500}{mL} of $0.9\%$ saline. (Hint: find solute mass, then use the van 't Hoff factor.)
\end{sikaoti}
